% ---------------------------------------------------------------------
% --- CAPÍTULO 4 - PROJETO E PROTÓTIPO
% ---
% --- Conteúdo extraído VERBATIM de TCC_PARKON_V3.pdf (páginas 20 a 30),
% --- fonte da verdade. Fragmento LaTeX (sem \documentclass/preâmbulo),
% --- incluído via \include em tese.tex. Arquivo em UTF-8.
% ---
% --- Título de origem (corpo do PDF): "4 PROJETO E PROTÓTIPO".
% ---   A classe PPGC_UFF formata o título do capítulo em maiúsculas; por isso
% ---   usa-se \chapter{Projeto e Protótipo} (mesma saída visível), mantendo a
% ---   coerência com os capítulos anteriores (Requisitos 1.2 e 1.6).
% ---   O label é \label{cap:resultados} (referenciado por \autoref{cap:resultados}
% ---   no cap1.tex) conforme definido na tarefa 4.4.
% ---
% --- ESTRUTURA DE SEÇÕES (hierarquia do PDF):
% ---   4   PROJETO E PROTÓTIPO        -> \chapter      (cap:resultados)
% ---   4.1 Criação do protótipo       -> \section      (sec:criacao-prototipo)
% ---   4.2 Criação do projeto         -> \section      (sec:criacao-projeto)
% ---   4.3 Telas do Aplicativo        -> \section      (sec:telas-aplicativo)
% ---   4.4 Build do Aplicativo        -> \section      (sec:build-aplicativo)
% ---
% --- MAPEAMENTO DE CITAÇÕES (fonte usa numeração [n]; convertidas para \cite):
% ---   [7] -> \cite{guarnieri2025}  GUARNIERI, P. L. Estudo de caso comparando o
% ---          MongoDB e o PostgreSQL... (seção 4.2)
% ---   [5] -> \cite{gradin2025}     GRADIN, F. A.; LAZZARETTI, A. T. Um estudo sobre
% ---          o desenvolvimento mobile utilizando Flutter (seção 4.4)
% ---   TODO: [CITATION] As referências da fonte NÃO trazem ano de publicação
% ---         explícito (apenas "Acesso em 08 de dez. 2025"). Os anos acima são
% ---         PROVISÓRIOS (acesso 2025) e devem ser confirmados/reconciliados nas
% ---         tarefas 7.1 e 7.2.
% ---
% --- FIGURAS (Figura 4 a 22): inseridos AMBIENTES PLACEHOLDER com \fbox e \label
% ---   para que os \autoref resolvam; as imagens reais serão extraídas/inseridas
% ---   na tarefa 5.1. Não há tabelas neste capítulo.
% ---
% --- REFERÊNCIAS CRUZADAS: todas as menções "Figura N" no corpo foram convertidas
% ---   para \autoref{fig:...} (Requisito 9.1). Como o template usa
% ---   \counterwithout{figure}{chapter}, a numeração é sequencial no documento;
% ---   o cap3.tex contém as Figuras 1-3, de modo que este capítulo produz
% ---   naturalmente "Figura 4" a "Figura 22", coincidindo com a numeração da fonte.
% ---
% --- ACRÔNIMOS (definidos em pre-textuais/acronimos.tex):
% ---   cli -> ocorre em 4.4 (primeira e única ocorrência da sigla no corpo do
% ---          documento). Na reconciliação de primeiro uso (tarefa 10.1) passou a
% ---          usar \acrfull{cli} (Requisitos 7.2 e 7.5), que renderiza
% ---          "Command Line Interface (CLI)"; removeu-se o parêntese literal
% ---          "(Command Line Interface)" para evitar duplicar a forma por extenso.
% ---   apk -> aparece apenas dentro de um bloco de código/comando verbatim
% ---          ("# Build para Android (APK)"); mantido literal (não convertido),
% ---          por ser conteúdo de linha de comando reproduzido VERBATIM.
% ---
% --- CÓDIGOS/COMANDOS: os blocos de comandos de build (Flutter/Go/Docker) são
% ---   reproduzidos VERBATIM via ambiente \begin{verbatim} (Requisito 1.1).
% ---------------------------------------------------------------------

\chapter{Projeto e Protótipo}
\label{cap:resultados}

\section{Criação do protótipo}
\label{sec:criacao-prototipo}

O protótipo do aplicativo Parkon foi desenvolvido com o objetivo de validar as principais funcionalidades propostas para a reserva de vagas de estacionamento em ambientes urbanos. Utilizando ferramentas de prototipagem, foram criadas telas que simulam a experiência do usuário, permitindo a visualização do fluxo de navegação, das opções de busca e reserva de vagas, bem como dos mecanismos de avaliação e pagamento. O protótipo serve como base para testes iniciais de usabilidade e para o refinamento dos requisitos do sistema, antes da implementação definitiva.

\begin{figure}[H]
   \centering
   \begin{minipage}[b]{0.48\linewidth}
      \centering
      
\includegraphics[height=0.34\textheight,keepaspectratio]{capitulos/figuras/tela-autenticacao.png}
      \caption{Tela de Autenticação}
      \label{fig:tela-autenticacao}
   \end{minipage}\hfill
   \begin{minipage}[b]{0.48\linewidth}
      \centering
      \includegraphics[height=0.34\textheight,keepaspectratio]{capitulos/figuras/tela-inicial.png}
      \caption{Tela Inicial do aplicativo}
      \label{fig:tela-inicial}
   \end{minipage}
\end{figure}

A \autoref{fig:tela-autenticacao} apresenta a tela de autenticação do aplicativo Parkon, onde o usuário pode realizar login ou cadastro, inclusive utilizando contas do Google ou Apple, em uma interface simples e moderna. Já a \autoref{fig:tela-inicial} mostra a tela inicial, composta por um mapa interativo e um campo de pesquisa ``Onde parar?'', permitindo ao usuário localizar estacionamentos próximos de forma prática e eficiente.

\begin{figure}[H]
   \centering
   \begin{minipage}[b]{0.48\linewidth}
      \centering
      \includegraphics[height=0.34\textheight,keepaspectratio]{capitulos/figuras/tela-pesquisa-vagas.png}
      \caption{Tela de Pesquisa de Vagas}
      \label{fig:tela-pesquisa-vagas}
   \end{minipage}\hfill
   \begin{minipage}[b]{0.48\linewidth}
      \centering
      \includegraphics[height=0.34\textheight,keepaspectratio]{capitulos/figuras/tela-selecao-estacionamento.png}
      \caption{Tela de Seleção de Estacionamento}
      \label{fig:tela-selecao-estacionamento}
   \end{minipage}
\end{figure}

A \autoref{fig:tela-pesquisa-vagas} apresenta a tela de pesquisa de vagas do aplicativo Parkon, onde o usuário pode inserir o destino desejado e visualizar uma lista de estacionamentos disponíveis na região, facilitando a escolha do local para estacionar. Já a \autoref{fig:tela-selecao-estacionamento} mostra a tela de seleção de estacionamento, exibindo detalhes como endereço, preço, avaliações e previsão de chegada, além de permitir ao usuário navegar até o estacionamento escolhido.

\begin{figure}[H]
   \centering
   \includegraphics[height=0.42\textheight,keepaspectratio]{capitulos/figuras/tela-configuracoes.png}
   \caption{Tela de Configurações}
   \label{fig:tela-configuracoes}
\end{figure}

\section{Criação do projeto}
\label{sec:criacao-projeto}

\subsection{Aplicativo Mobile (Flutter)}
\label{sec:criacao-projeto-mobile}

O desenvolvimento do projeto do aplicativo Parkon foi realizado utilizando o framework Flutter, escolhido por sua capacidade de criar aplicações móveis multiplataforma com uma única base de código, facilitando a manutenção e a escalabilidade do sistema. Durante a implementação, um dos principais desafios enfrentados foi a integração do serviço de geolocalização do Google, essencial para identificar a localização do usuário e das vagas disponíveis em tempo real. Para garantir o desempenho e a disponibilidade das informações, foi adotado um mecanismo de cache utilizando o banco de dados MongoDB~\cite{guarnieri2025}, permitindo o armazenamento temporário dos dados de localização e otimizando o acesso mesmo em situações de instabilidade de conexão. Essa abordagem proporcionou maior eficiência ao aplicativo, reduzindo o tempo de resposta e melhorando a experiência do usuário.

\subsection{Portal Web Administrativo (Next.js)}
\label{sec:criacao-projeto-web}

O portal web administrativo foi desenvolvido utilizando React com o framework Next.js, permitindo que gestores de estacionamento realizem operações de cadastro, configuração e acompanhamento de reservas por meio de uma interface web responsiva.

A escolha do Next.js se deu pela facilidade de integração com a \acrshort{api} REST já existente no backend em Go, além de oferecer renderização do lado do servidor, o que proporciona melhor desempenho no carregamento inicial das páginas. A interface foi construída com componentes reutilizáveis, seguindo padrões de design responsivo para garantir acessibilidade em diferentes dispositivos.

O portal compartilha a mesma infraestrutura de backend e banco de dados (MongoDB) utilizada pelo aplicativo mobile, garantindo consistência dos dados entre ambas as plataformas. A autenticação é realizada por meio do mesmo serviço de autenticação, permitindo que as credenciais do administrador sejam utilizadas de forma unificada.

\section{Telas do Aplicativo}
\label{sec:telas-aplicativo}

\subsection{Telas do Aplicativo Mobile}
\label{sec:telas-mobile}

\begin{figure}[H]
   \centering
   \begin{minipage}[b]{0.48\linewidth}
      \centering
      \includegraphics[height=0.34\textheight,keepaspectratio]{capitulos/figuras/emulador-login.png}
      \caption{Aplicação rodando no emulador - Login}
      \label{fig:emulador-login}
   \end{minipage}\hfill
   \begin{minipage}[b]{0.48\linewidth}
      \centering
      \includegraphics[height=0.34\textheight,keepaspectratio]{capitulos/figuras/emulador-cadastro.png}
      \caption{Aplicação rodando no emulador - Cadastro}
      \label{fig:emulador-cadastro}
   \end{minipage}
\end{figure}

A \autoref{fig:emulador-login} mostra a tela de login do aplicativo Parkon em funcionamento no emulador, destacando o campo para inserção de e-mail e opções de autenticação, incluindo Google e Apple, em um layout moderno e funcional. A \autoref{fig:emulador-cadastro} apresenta a tela de cadastro, onde o usuário informa nome e senha para criar sua conta.

\begin{figure}[H]
   \centering
   
\includegraphics[height=0.34\textheight,keepaspectratio]{capitulos/figuras/emulador-confirmacao-email.png}
   \caption{Aplicação rodando no emulador - Confirmação de e-mail}
   \label{fig:emulador-confirmacao-email}

   \vspace{0.6cm}

   \includegraphics[width=0.85\linewidth]{capitulos/figuras/emulador-email-codigo.png}
   \caption{Aplicação rodando no emulador – E-mail com o código}
   \label{fig:emulador-email-codigo}
\end{figure}

A \autoref{fig:emulador-confirmacao-email} apresenta a tela de confirmação de e-mail do aplicativo Parkon, onde o usuário deve inserir o código recebido para validar o cadastro. Já a \autoref{fig:emulador-email-codigo} mostra o e-mail enviado ao usuário com o código de verificação necessário para concluir o processo de registro.

\begin{figure}[H]
   \centering
   \begin{minipage}[b]{0.48\linewidth}
      \centering
      \includegraphics[height=0.34\textheight,keepaspectratio]{capitulos/figuras/emulador-sucesso-cadastro.png}
      \caption{Aplicação rodando no emulador - Sucesso no cadastro}
      \label{fig:emulador-sucesso-cadastro}
   \end{minipage}\hfill
   \begin{minipage}[b]{0.48\linewidth}
      \centering
      \includegraphics[height=0.34\textheight,keepaspectratio]{capitulos/figuras/emulador-tela-inicial.png}
      \caption{Aplicação rodando no emulador – Tela ínicial}
      \label{fig:emulador-tela-inicial}
   \end{minipage}
\end{figure}

A \autoref{fig:emulador-sucesso-cadastro} mostra a tela de confirmação de sucesso no cadastro do usuário, indicando que o processo foi concluído corretamente no aplicativo Parkon. Já a \autoref{fig:emulador-tela-inicial} apresenta a tela inicial do aplicativo após o login, exibindo o mapa da região e o campo de pesquisa para localização de vagas.

\begin{figure}[H]
   \centering
   \begin{minipage}[b]{0.48\linewidth}
      \centering
      \includegraphics[height=0.34\textheight,keepaspectratio]{capitulos/figuras/emulador-configuracoes.png}
      \caption{Aplicação rodando no emulador - Configurações}
      \label{fig:emulador-configuracoes}
   \end{minipage}\hfill
   \begin{minipage}[b]{0.48\linewidth}
      \centering
      
\includegraphics[height=0.34\textheight,keepaspectratio]{capitulos/figuras/emulador-perfil.png}
      \caption{Aplicação rodando no emulador – Perfil do usuário}
      \label{fig:emulador-perfil}
   \end{minipage}
\end{figure}

A \autoref{fig:emulador-configuracoes} apresenta a tela de configurações do aplicativo Parkon, onde o usuário pode acessar opções como perfil, privacidade, histórico de buscas, pagamentos e ajustes gerais. A \autoref{fig:emulador-perfil} mostra a tela de perfil do usuário, exibindo dados pessoais e a opção de encerrar a conta.

\begin{figure}[H]
   \centering
   \begin{minipage}[b]{0.48\linewidth}
      \centering
      \includegraphics[height=0.34\textheight,keepaspectratio]{capitulos/figuras/emulador-pesquisa.png}
      \caption{Aplicação rodando no emulador - Pesquisa de estacionamento}
      \label{fig:emulador-pesquisa}
   \end{minipage}\hfill
   \begin{minipage}[b]{0.48\linewidth}
      \centering
      \includegraphics[height=0.34\textheight,keepaspectratio]{capitulos/figuras/emulador-lista.png}
      \caption{Aplicação rodando no emulador - Lista de estacionamentos}
      \label{fig:emulador-lista}
   \end{minipage}
\end{figure}

A \autoref{fig:emulador-pesquisa} apresenta a tela de pesquisa de estacionamento do aplicativo Parkon, onde o usuário pode digitar o endereço desejado para localizar vagas disponíveis na região. A \autoref{fig:emulador-lista} mostra a lista de estacionamentos encontrados, exibindo opções próximas ao destino informado pelo usuário.

\begin{figure}[H]
   \centering
   \begin{minipage}[b]{0.48\linewidth}
      \centering
      \includegraphics[height=0.34\textheight,keepaspectratio]{capitulos/figuras/emulador-detalhes.png}
      \caption{Aplicação rodando no emulador - Detalhes do estacionamento}
      \label{fig:emulador-detalhes}
   \end{minipage}\hfill
   \begin{minipage}[b]{0.48\linewidth}
      \centering
      \includegraphics[height=0.34\textheight,keepaspectratio]{capitulos/figuras/emulador-comentarios.png}
      \caption{Aplicação rodando no emulador - Comentários}
      \label{fig:emulador-comentarios}
   \end{minipage}
\end{figure}

A \autoref{fig:emulador-detalhes} mostra a tela de detalhes do estacionamento no aplicativo Parkon, exibindo informações como endereço, distância, avaliações e opções para ver comentários ou reportar problemas. A \autoref{fig:emulador-comentarios} apresenta a tela de comentários, onde o usuário pode visualizar e adicionar opiniões sobre o estacionamento selecionado.

\begin{figure}[H]
   \centering
   \begin{minipage}[b]{0.48\linewidth}
      \centering
      \includegraphics[height=0.34\textheight,keepaspectratio]{capitulos/figuras/emulador-reportar.png}
      \caption{Aplicação rodando no emulador – Reportar problemas}
      \label{fig:emulador-reportar}
   \end{minipage}\hfill
   \begin{minipage}[b]{0.48\linewidth}
      \centering
      \includegraphics[height=0.34\textheight,keepaspectratio]{capitulos/figuras/emulador-navegacao.png}
      \caption{Aplicação rodando no emulador - Tela de navegação do google maps}
      \label{fig:emulador-navegacao}
   \end{minipage}
\end{figure}

A \autoref{fig:emulador-reportar} apresenta a tela de reportar problemas do aplicativo Parkon, permitindo ao usuário informar à equipe sobre dificuldades encontradas na rota até o estacionamento. A \autoref{fig:emulador-navegacao} mostra a tela de navegação do Google Maps integrada ao aplicativo, exibindo o trajeto detalhado até o estacionamento selecionado.

\begin{figure}[H]
   \centering
   \begin{minipage}[b]{0.48\linewidth}
      \centering
      \includegraphics[height=0.34\textheight,keepaspectratio]{capitulos/figuras/app-reservar.jpeg}
      \caption{Aplicação rodando no emulador - Detalhes do estacionamento credenciado}
      \label{fig:app-reservar}
   \end{minipage}\hfill
   \begin{minipage}[b]{0.48\linewidth}
      \centering
      \includegraphics[height=0.34\textheight,keepaspectratio]{capitulos/figuras/app-reservar-data.jpeg}
      \caption{Aplicação rodando no emulador - Seleção de data e hora da reserva}
      \label{fig:app-reservar-data}
   \end{minipage}
\end{figure}

\begin{figure}[H]
   \centering
   \includegraphics[height=0.34\textheight,keepaspectratio]{capitulos/figuras/app-reservado.jpeg}
   \caption{Aplicação rodando no emulador - Reserva solicitada aguardando confirmação}
   \label{fig:app-reservado}
\end{figure}

A \autoref{fig:app-reservar} apresenta a tela de detalhes de um estacionamento credenciado na plataforma, exibindo a localização no mapa, endereço, distância, avaliação e os botões de ação ``Ver Comentários'', ``Reportar Problema'', ``Navegar'' e ``Reservar''. A \autoref{fig:app-reservar-data} mostra a tela de reserva, onde o usuário seleciona a data e hora de início e fim da reserva e confirma com o botão ``Solicitar Reserva''. A \autoref{fig:app-reservado} apresenta a confirmação da solicitação, com o banner informando que a reserva foi solicitada e aguarda confirmação do administrador do estacionamento.

\subsection{Telas do Portal Web}
\label{sec:telas-web}

\begin{figure}[H]
   \centering
   \includegraphics[width=0.85\linewidth]{capitulos/figuras/portal-login.png}
   \caption{Portal Web - Tela de Login}
   \label{fig:portal-login}
\end{figure}

\begin{figure}[H]
   \centering
   \includegraphics[width=0.85\linewidth]{capitulos/figuras/portal-cadastro.png}
   \caption{Portal Web - Tela de Cadastro}
   \label{fig:portal-cadastro}
\end{figure}

A \autoref{fig:portal-login} apresenta a tela de login do portal web administrativo ParkOn, onde o administrador do estacionamento pode inserir e-mail e senha para acessar o painel de gestão. Já a \autoref{fig:portal-cadastro} mostra a tela de cadastro, onde o administrador informa nome completo, e-mail, telefone, CPF, senha e confirmação de senha para criar sua conta na plataforma.

\begin{figure}[H]
   \centering
   \includegraphics[width=0.85\linewidth]{capitulos/figuras/portal-confirmacao-email.png}
   \caption{Portal Web - E-mail com código de verificação}
   \label{fig:portal-confirmacao-email}
\end{figure}

\begin{figure}[H]
   \centering
   \includegraphics[width=0.85\linewidth]{capitulos/figuras/portal-confirmacao.png}
   \caption{Portal Web - Modal de confirmação de código}
   \label{fig:portal-confirmacao}
\end{figure}

A \autoref{fig:portal-confirmacao-email} mostra o e-mail enviado ao administrador com o código de verificação necessário para concluir o cadastro na plataforma. A \autoref{fig:portal-confirmacao} apresenta o modal de confirmação exibido sobre a tela de cadastro, onde o administrador insere o código recebido por e-mail para validar sua conta.

\begin{figure}[H]
   \centering
   \includegraphics[width=0.85\linewidth]{capitulos/figuras/portal-setup-1.png}
   \caption{Portal Web - Configuração inicial (Bem Vindo)}
   \label{fig:portal-setup-1}
\end{figure}

\begin{figure}[H]
   \centering
   \includegraphics[width=0.85\linewidth]{capitulos/figuras/portal-setup-2.png}
   \caption{Portal Web - Configuração inicial (Dados da Empresa)}
   \label{fig:portal-setup-2}
\end{figure}

\begin{figure}[H]
   \centering
   \includegraphics[width=0.85\linewidth]{capitulos/figuras/portal-setup-3.png}
   \caption{Portal Web - Configuração inicial (Primeiro Estacionamento)}
   \label{fig:portal-setup-3}
\end{figure}

Após a validação do cadastro, o administrador é direcionado para o assistente de configuração inicial. A \autoref{fig:portal-setup-1} apresenta a tela de boas-vindas com o botão ``Iniciar''. A \autoref{fig:portal-setup-2} mostra o formulário de cadastro da empresa, com campos para razão social, CNPJ e telefone comercial, além de uma declaração de responsabilidade sobre os dados informados. A \autoref{fig:portal-setup-3} apresenta o cadastro do primeiro estacionamento, onde o administrador informa o nome do estabelecimento e o endereço, com integração ao Google Maps para posicionamento no mapa.

\begin{figure}[H]
   \centering
   \includegraphics[width=0.85\linewidth]{capitulos/figuras/portal-dashboard.png}
   \caption{Portal Web - Dashboard}
   \label{fig:portal-dashboard}
\end{figure}

A \autoref{fig:portal-dashboard} apresenta o painel principal do portal administrativo, exibindo métricas como total de estacionamentos, total de vagas, taxa de ocupação e reservas pendentes, além de um gráfico de reservas por hora. O menu lateral permite acesso rápido às demais funcionalidades do sistema.

\begin{figure}[H]
   \centering
   \includegraphics[width=0.85\linewidth]{capitulos/figuras/portal-empresas.png}
   \caption{Portal Web - Gerenciamento de Empresas}
   \label{fig:portal-empresas}
\end{figure}

A \autoref{fig:portal-empresas} mostra a tela de gerenciamento de empresas, onde o administrador pode visualizar as empresas cadastradas com informações como nome, CNPJ, telefone, status de verificação e ações disponíveis, além de criar novas empresas.

\begin{figure}[H]
   \centering
   \includegraphics[width=0.85\linewidth]{capitulos/figuras/portal-estacionamentos.png}
   \caption{Portal Web - Lista de Estacionamentos}
   \label{fig:portal-estacionamentos}
\end{figure}

A \autoref{fig:portal-estacionamentos} apresenta a lista de estacionamentos cadastrados, exibindo nome, endereço, vagas totais, vagas ocupadas e status de verificação, com a opção de criar novos estacionamentos.

\begin{figure}[H]
   \centering
   \includegraphics[width=0.85\linewidth]{capitulos/figuras/portal-edicao-estacionamento.png}
   \caption{Portal Web - Detalhes do Estacionamento (Informações)}
   \label{fig:portal-edicao-estacionamento}
\end{figure}

\begin{figure}[H]
   \centering
   \includegraphics[width=0.85\linewidth]{capitulos/figuras/portal-edicao-estacionamento-horario.png}
   \caption{Portal Web - Detalhes do Estacionamento (Horários de Funcionamento)}
   \label{fig:portal-edicao-estacionamento-horario}
\end{figure}

A \autoref{fig:portal-edicao-estacionamento} mostra a tela de detalhes do estacionamento com informações gerais (nome, endereço, status) e a localização no mapa. A \autoref{fig:portal-edicao-estacionamento-horario} apresenta a configuração dos horários de funcionamento por dia da semana, com opção de marcar dias como fechado.

\begin{figure}[H]
   \centering
   \includegraphics[width=0.85\linewidth]{capitulos/figuras/portal-edicao-estacionamento-vagas.png}
   \caption{Portal Web - Gerenciamento de Vagas}
   \label{fig:portal-edicao-estacionamento-vagas}
\end{figure}

A \autoref{fig:portal-edicao-estacionamento-vagas} apresenta a tela de gerenciamento de vagas do estacionamento, onde o administrador pode visualizar o identificador e o status de cada vaga, alterar sua disponibilidade e adicionar novas vagas.

\begin{figure}[H]
   \centering
   \includegraphics[width=0.85\linewidth]{capitulos/figuras/portal-equipe.png}
   \caption{Portal Web - Gerenciamento de Equipe}
   \label{fig:portal-equipe}
\end{figure}

A \autoref{fig:portal-equipe} mostra a tela de gerenciamento de equipe, onde o administrador visualiza os membros com seus respectivos papéis e status, além de poder convidar novos colaboradores para auxiliar na gestão do estacionamento.

\begin{figure}[H]
   \centering
   \includegraphics[width=0.85\linewidth]{capitulos/figuras/portal-reservas.png}
   \caption{Portal Web - Gerenciamento de Reservas}
   \label{fig:portal-reservas}
\end{figure}

A \autoref{fig:portal-reservas} apresenta a tela de gerenciamento de reservas, exibindo informações como e-mail do usuário, estacionamento, vaga, período, status e ações disponíveis, permitindo ao administrador acompanhar e gerenciar as reservas realizadas pelos motoristas.

\section{Build do Aplicativo}
\label{sec:build-aplicativo}

O processo de build da plataforma Parkon envolve três componentes distintos: o aplicativo mobile (Flutter), o backend (Go) e o portal web administrativo (Next.js). Para cada um, foram utilizados comandos \acrfull{cli} que automatizam a criação dos artefatos para distribuição e execução.

\subsection{Build do Aplicativo Mobile (Flutter)}
\label{sec:build-mobile}

O aplicativo mobile foi desenvolvido utilizando o framework Flutter~\cite{gradin2025}, que permite a geração de pacotes para múltiplas plataformas, como Android e iOS, a partir de uma única base de código. Exemplos de comandos utilizados:

\begin{verbatim}
# Build para Android (APK)
flutter build apk --release

# Build para Android (App Bundle)
flutter build appbundle --release

# Build para iOS (necessário macOS)
flutter build ios --release
\end{verbatim}

Esses comandos geram os pacotes necessários para publicação nas lojas de aplicativos ou para testes em dispositivos físicos.

\subsection{Build do Backend (Go)}
\label{sec:build-backend}

O backend do sistema foi desenvolvido em Go, proporcionando alta performance e facilidade de escalabilidade. O processo de build do backend é realizado via linha de comando:

\begin{verbatim}
# Compilar o backend
go build -o parkon-backend main.go

# Executar o backend localmente
./parkon-backend

# Para rodar em ambiente de produção com Docker:
docker build -t parkon-backend .
docker run -d -p 8080:8080 parkon-backend
\end{verbatim}

Esses comandos garantem que o backend esteja pronto para execução local ou implantação em ambientes cloud, facilitando a integração com o aplicativo mobile e o portal web.

\subsection{Build do Portal Web (Next.js)}
\label{sec:build-portal-web}

O processo de build do portal web é realizado via linha de comando:

\begin{verbatim}
# Instalar dependências
npm install

# Executar em ambiente de desenvolvimento
npm run dev

# Build para produção
npm run build

# Iniciar em produção
npm start
\end{verbatim}
